\documentclass[12pt]{article}
\usepackage[utf8]{inputenc}
\usepackage[T1]{fontenc}
\usepackage{hyperref}
\title{Swati Language Resources for K-12 Curriculum Development}
\author{Compiled by Claude Sonnet 5.5 (claude-sonnet-5-5) for Andr\'as Kornai}
\date{}
\begin{document}
\maketitle

\section{Introduction}
Swati (siSwati, ISO 639-3: ssw) is a Bantu language and an official national language of Eswatini. It is also spoken in South Africa. Search results report that Eswatini moved in 2011 to use siSwati as the medium of instruction in the early grades and as a core subject through senior secondary.

\section{Primary Educational Resources}
\begin{itemize}
\item Oxford Kukhanyo, siSwati Home Language, Grades 1--3: learner's books and teacher's guides. It is listed for South Africa's CAPS curriculum, so its fit for Eswatini is not confirmed. \url{https://oxford.co.za/schools/oxford-caps-approved/foundation-phase/siswati-home-language}
\item Oxford First Bilingual Dictionary: Siswati and English, aimed at learners from Grade 2. \url{https://oxford.co.za/shop/resources-and-home-learning/resource-book/9780195992915-oxford-first-bilingual-dictionary-siswati-and-english/}
\end{itemize}

\section{Linguistic and Reference Materials}
\begin{itemize}
\item Cambridge Language Centre page for \emph{Asikhulume SiSwati} (Dlamini, 2010): a beginner course with 17 lessons and a two-way glossary. \url{https://www.langcen.cam.ac.uk/resources/langs/data/ss_102.html}
\item African Union library: ``Raising the profile of siSwati as a national language''. \url{https://library.au.int/raising-profile-siswati-national-language}
\item Journal of Languages and Language Teaching (Jan 2023), a study of siSwati as a first-language subject. \url{https://doaj.org/article/b9f5b39c07ae4491a838efe541150b53}
\end{itemize}

\section{Religious and Community Resources}
None were found in the search.

\section{Digital and Audio Resources}
\begin{itemize}
\item Peace Corps SiSwati language course, listed on fsi-language-courses.org. It sits on a commercial site, so check its provenance before relying on it. \url{https://www.fsi-language-courses.org/siswati/courses/peace-corps-siswati-language-course}
\end{itemize}

\section{Curriculum Development Context}
siSwati has official status and is used in teaching, and commercial and academic materials exist. Most of the textbooks found target South Africa. Eswatini-specific materials, such as those from the Ministry of Education or the SiSwati Language Board, should be sought directly. The search did not confirm any of them.

\end{document}